% Figure: slope field for the logistic equation, with two integral
% curves shown to demonstrate convergence to the carrying capacity.
\begin{figure}[htbp]
\centering
\begin{tikzpicture}[
  axis/.style={draw=engblack, -{Stealth}, thick},
  field/.style={draw=enggray, thin, -{Latex[length=1.4mm, width=1.0mm]}},
  curveA/.style={draw=engblue, very thick},
  curveB/.style={draw=engred, very thick},
  curveC/.style={draw=engamber, very thick},
  label/.style={font=\small},
]
% Axes (scaled: x from 0 to 10 (time), y from 0 to 6 (N/K, where K=1 -> 4 units))
\draw[axis] (0,0) -- (11,0) node[anchor=west, label] {$t$};
\draw[axis] (0,0) -- (0,5.4) node[anchor=south, label] {$N$};

% Ticks
\foreach \x in {0, 2, 4, 6, 8, 10} {
  \draw (\x,-0.07) -- (\x,0.07);
  \node[font=\scriptsize, below] at (\x,-0.07) {\x};
}
\foreach \y/\h in {0/0, 1/0.25, 2/0.50, 3/0.75, 4/1.00} {
  \draw (-0.07,\y) -- (0.07,\y);
  \node[font=\scriptsize, left] at (-0.07,\y) {\h};
}

% Carrying capacity line K=1 (at y=4)
\draw[draw=enggray, dashed, thin] (0,4) -- (11,4);
\node[font=\scriptsize, enggray, anchor=west] at (10.5, 4.2) {$K$};

% Slope field: short arrows; dN/dt = N(1 - N/4) (r=1, K=4 in plot units)
% Plot units: y in 0..4 is N in 0..1
\foreach \tx in {0.4, 1.2, 2.0, 2.8, 3.6, 4.4, 5.2, 6.0, 6.8, 7.6, 8.4, 9.2, 10.0} {
  \foreach \ny in {0.2, 0.6, 1.0, 1.4, 1.8, 2.2, 2.6, 3.0, 3.4, 3.8, 4.2, 4.6, 5.0} {
    % slope = N(1 - N/4) where N = ny/4 maps to actual N
    % Use ny as the rendered y; slope_render = (ny/4)*(1 - ny/4)
    \pgfmathsetmacro{\slp}{(\ny/4)*(1 - \ny/4)}
    \pgfmathsetmacro{\norm}{sqrt(1 + (\slp)^2)}
    \pgfmathsetmacro{\dx}{0.35/\norm}
    \pgfmathsetmacro{\dy}{0.35*\slp/\norm}
    \draw[field] (\tx - \dx, \ny - \dy) -- (\tx + \dx, \ny + \dy);
  }
}

% Three integral curves: logistic solutions with K=1, r=1, three N0
% Plot uses N_plot = 4 * K * N0 * exp(t) / (K + N0 (exp(t) - 1))
% A: N0 = 0.05; B: N0 = 0.5; C: N0 = 1.5 (overshoot)
\draw[curveA] plot[domain=0:6, samples=80]
  ({\x}, {4 / (1 + ((1 - 0.05)/0.05) * exp(-\x))});
\draw[curveB] plot[domain=0:6, samples=80]
  ({\x}, {4 / (1 + ((1 - 0.5)/0.5) * exp(-\x))});
\draw[curveC] plot[domain=0:6, samples=80]
  ({\x}, {4 / (1 + ((1 - 1.5)/1.5) * exp(-\x))});

% Labels
\node[engblue, font=\scriptsize, anchor=west] at (4.5, 3.0) {$N_{0}=0.05\,K$};
\node[engred, font=\scriptsize, anchor=west] at (3.2, 4.45) {$N_{0}=0.5\,K$};
\node[engamber, font=\scriptsize, anchor=west] at (1.0, 5.1) {$N_{0}=1.5\,K$ (over)};

\end{tikzpicture}
\caption[Slope field for the logistic equation]{Slope field for the
logistic equation $\dot{N}=rN(1-N/K)$ with $r=1$ and $K=1$, with
three integral curves drawn at initial conditions $N_{0}=0.05\,K$,
$N_{0}=0.5\,K$, and $N_{0}=1.5\,K$. All solutions converge to the
carrying capacity $K$. The inflection point at $N=K/2$ marks the
maximum growth rate; trajectories below approach the inflection
from below and accelerate, trajectories above decay monotonically.
The slope field summarises every solution at once; the integral
curves are particular solutions selected by initial condition.}
\label{fig:vol02:ch07:slope-field}
\end{figure}
